\section{Four polynomial links between the lattice and the Case III seed}

Fix an original Case III datum $D$, an acceptable, tame tower $T_0$
with the source four-factor property, an embedding $\iota$, the ambient
residue field $k$, and a global residual pair $c_0$ in the prescribed
lattice shape. Assume the cohomological input, the fitness implication
and the seed hypothesis $\mathrm{SeedIVAlt}$. Our aim is one extension
$T_1$ of $T_0$ with a property for \emph{every} later deep tower $T$:
on $T$ one integral seed must work for every admissible restriction of
$c_0$. The seed hypothesis gives a residual pair before $T_1$ is
chosen. The integral seed will only be chosen on $T$.

Use zero-based matrix indices. Write $\alpha,\beta,\sigma,\tau$ for
the source weights $\mathrm{wAl},\mathrm{wBe},\mathrm{wPs},\mathrm{wTh}$.
The support sets for the two types of lines and their intersection are
\[
 N_P=\{\beta,-\alpha,-\tau,-\sigma\},\qquad
 N_M=\{-\alpha,-\beta,-\sigma,-\tau\},\qquad
 Q_0=\{-\alpha,-\sigma,-\tau\}.
\]
Support restricts the off-diagonal entries of the representation.
Saying a pair is in $Q_0$ also restricts its local lower unitriangular
flag matrices. Write $E_\lambda$ for the existence of a nonzero
representation entry of weight $\lambda$. Two pairs have
$\mathrm{SamePs}$ if the $20$ entries agree both on the group and in
each local flag. $\mathrm{SameAl}$ means the corresponding equality of
$10$ entries.

% BEGIN RFCR1
\paragraph{RFCR1. Preserving one simple root while introducing the other.}
On the fit initial tower apply the seed hypothesis at the given
$\iota,k$ (with the original ambient datum). It supplies a split prime
$w$ and a residual pair $s$ supported on $N_M$ with auxiliary set
$\{w\}$. We construct $p_2$ on a larger finite split set $W_1$, in
$Q_0$, with $\mathrm{SameAl}(p_2,s)$ and $E_{-\sigma}(p_2)$.

Put $x(g)=s.r(g)_{10}/\bar\chi_0(g)$ and $m_v=s.n_{v,10}$, where
$\bar\chi_i$ are the actual residual diagonal characters. Lower
triangular multiplication makes $x$ a cocycle for
$\rho=\bar\chi_1/\bar\chi_0$. Normalization gives $x(g_0)=0$.
The first column of the local ordinary flag gives
$x(h)=(1-\rho(h))m_v$ on the decomposition group. We need an
independent cocycle $y$ for $\nu=\bar\chi_2/\bar\chi_0$.

Here is the interchange of roots in the cup construction of UX3.
Let $k_0$ be the image of the finite coefficient residue field, $L$ the
finite $k_0$-span of the values of $x$. This span is finite-dimensional
by compactness and local constancy. On $M=k_0\oplus L$ let
\[
 g(a,b)=\bigl(\nu(g)a,\ \rho(g)\nu(g)b+\nu(g)a\,x(g)\bigr).
\]
The cocycle rule for $x$ gives the action law. Use the finite-module
argument (UX.15) for $M$ and $(1,0)$. At the old places impose the
zero local class. At each added splitting prime allow the additional
one-dimensional tame inertia image spanned by $(1,0)$. Adding enough
primes makes the lower bound for the Selmer cardinality exceed the
number of cocycles unramified at all the added primes. The resulting
cocycle $(y,z)$ takes a nonzero multiple of $(1,0)$ on some such
inertia, where both characters are trivial. Subtract the coboundary
\[
 g\longmapsto\bigl((1-\nu(g))b,\
 (1-\rho(g)\nu(g))c-\nu(g)b\,x(g)\bigr)
\]
with $b=y(g_0)/(1-\nu(g_0))$ and
$c=z(g_0)/(1-\rho(g_0)\nu(g_0))$. Standardness separates both
roots at $g_0$. This sets both coordinates to zero there and preserves
the witnessed nonzero $y$. Write $n_v$ for its resulting local
constant, so $y(h)=(1-\nu(h))n_v$.

The cross term of $z(gh)$ is $x(g)\nu(g)y(h)$. To put the coordinates
back in the original matrix order set $z'=xy-z$. Expanding
$x(gh)y(gh)$ gives cross term $y(g)\rho(g)x(h)$ for $z'$.
The same subtraction at a local group replaces its central constant
$o_v$ by $o'_v=m_v n_v-o_v$. Use the matrices
\[
 U(x,y,t)=\begin{pmatrix}1&0&0&0\\ x&1&0&0\\ y&0&1&0\\ t&y&-x&1\end{pmatrix},
 \qquad r(g)=U(x,y,2z'-xy)\operatorname{diag}(\bar\chi_i(g))
\]
and the local matrices $U(m_v,n_v,2o'_v-m_vn_v)$. Their product
adds $yx'-xy'$ to the two central entries, exactly the cup identity
just obtained. The local identity gives the prescribed diagonal
characters after conjugation. The matrices are symplectic up to the
given diagonal multiplier, normalized at $g_0$, and trivial on the
required forbidden inertia. They define the promised $p_2$: the
$10$ data are unchanged and $r_{20}=y\bar\chi_0$ is nonzero somewhere.
% END RFCR1

% BEGIN RFCR2
\paragraph{RFCR2. Two simultaneous modifications before choosing the later tower.}
First restrict $c_0$ to $T_0$ and move the global ordinary flag to
the chosen places over $\ell$. More explicitly, if $g$ moves the
chosen place to a place above the new one, factor $c_0.r(g)c_0.n$
into a lower unitriangular matrix and an upper matrix. The successive
leading minors are units by the lattice support, as in UX2 and UX11.
Use the lower factor as flag; the upper factor conjugates upper matrices
to upper matrices and preserves their diagonal characters. Restriction
preserves the multiplier and continuity. The chosen $g_0$ is the
initial separating element. The ramification, inertia and diagonal
conditions of the restricted frame follow from those of the lattice
and fitness. This supplies $p_0$ over $c_0$, residual on the empty
auxiliary set, with representation supported on $N_P$ and flags on
$Q_0$. In particular the lower factor has zero $21$ entry (UX.13).

Regard $p_0,p_2$ on $W_1$ by weakening the forbidden-inertia test.
Each is supported on $N_P$ with flags on $Q_0$. For both take the
$20$ entry divided by $\bar\chi_0$ and its $20$ flag coordinate.
These are normalized $-\sigma$ cocycles and their local constants
by the first-column identity. Apply the \emph{simultaneous}
construction (UX.14)--(UX.19) to the two coordinates. It produces
one $W_2\supset W_1$, one nonzero normalized $-\alpha$ cocycle
and its local constants, and two $Q_0$ pairs $p_3,p_4$. Thus
\[
 \mathrm{SamePs}(p_0,p_3),\qquad
 \mathrm{SameAl}(p_3,p_4),\qquad
 \mathrm{SamePs}(p_4,p_2),\qquad E_{-\alpha}(p_3).
\]
In particular equality of the middle $10$ entries includes their local
constants. Choosing the two modifications independently would not
give this equality. All primes of $W_2$ split in the precise sense
of $T_0$.

Choose $T_1$ now. For each $w'\in W_2$ adjoin to the totally real
tower field the abelian degree-$\ell$ characters unramified outside
$w'$ by taking the compositum of their fixed fields of kernels.
Here $w'\ne\ell$. At $\ell$ the rational inertia group acts
nontrivially on $\mu_\ell$, whereas splitting requires the
Teichm\"uller cyclotomic character to be trivial there.
The character fields lie in the cyclic degree-at-most-$\ell$ part
of $\mathbb Q(\mu_{w'})$ by reciprocity. Their compositum over the
initial field is totally real, of exponent dividing $\ell$ and
unramified over $\ell$. Use the same label and auxiliary base field
and the $\ell$th power of the marked element in the new tower
(the tower extension construction in UX1). These operations require
acceptability and tameness of the initial tower and the specified
splitting primes, all present here. The tower remains tame at $\ell$;
its field kills every permitted degree-$\ell$ character at every
$w'\in W_2$. This is the source condition $\mathrm{Dies}$, also
valid at all later towers.
% END RFCR2

% BEGIN RFCR3
\paragraph{RFCR3. Restriction preserves the two matched coordinates and their witnesses.}
Fix $T\ge T_1$. The tower relation is transitive and preserves
fitness and acceptability. Restrict each $p_i$ for $i=0,2,3,4$ to
the new group and use the lower factors of its moved flags as in
RFCR2. Call the resulting pairs $q_i$. For the $N_M$ shape the
matrix product of a lower triangular representation and the old
lower flag is already lower triangular with invertible diagonal.
For the $N_P$ shape the conjugated flag has its lower factor in
$Q_0$ by the same minor calculation. Changes of the element moving
a place only multiply the moved flag on the right by an upper
matrix, so its normalized lower factor is unique. The original
ambient residual field and the four characters stay fixed.

At each auxiliary prime all diagonal characters on inertia are $1$.
The two allowed shapes are triangular in their respective orderings,
so $(r(g)-1)^4=0$ on inertia. Here $\ell\ge5$ by the four-factor
hypothesis, hence $r(g)^\ell=1$. The image is finite by local
constancy and compactness. It is an $\ell$-group (each of its elements
has order dividing $\ell$), hence kills wild inertia and factors
through the cyclic order-$\ell$ quotient of tame inertia. Splitting
gives $w'\equiv1\bmod\ell$ and identifies the local base field with
$\mathbb Q_{w'}$; the cyclotomic degree-$\ell$ character at $w'$
is nontrivial on this quotient. The $\mathrm{Dies}$ condition thus
kills the residual inertia on $T$. Away from $W_2$ the old forbidden
inertia was already killed. The new marked element restricts to an
$\ell$-power of the old one, so is still diagonal; the new local
characters are those at the original places by conjugacy. This gives
the residual conditions with the empty auxiliary set.

For an $N_P$ pair the first column of the factorization gives
\[
 n'_{20}r_{00}(g)=r_{20}(g)+r_{22}(g)n_{20}.
\]
For an $N_M$ pair one similarly has
$n'_{10}r_{00}(g)=r_{10}(g)+r_{11}(g)n_{10}$.
Both identities use the actual place mover. Since the diagonal entries
are the common nonzero characters, the first identity transports
$\mathrm{SamePs}$ and the second transports $\mathrm{SameAl}$.
A second restriction of the same global pair also has exactly the
same lower flags, by uniqueness of those factors. Consequently any
residual $u$ above $c_0$ with the stipulated flag compatibility is $q_0$.

We record why the needed nonzero entries survive. On the indicated
supports, the entry $d=r_{20}$ for $N_P$, or $d=r_{10}$ for $N_M$,
satisfies
$d(gh)=a(g)d(h)+d(g)b(h)$ with its two diagonal characters $a,b$.
Take $x$ with $d(x)\ne0$, and $y$ with $a(y)\ne b(y)$ by separation.
If $a(x)=b(x)$ and $d(y)=0$, use $z=yx$; if $d(y)\ne0$ in this
case use $z=y$. Otherwise take $z=x$. In each case $d(z)\ne0$ and
$a(z)\ne b(z)$. For $N=\ell^m$ the power recurrence gives
\[
 d(z^N)\bigl(a(z)-b(z)\bigr)
       =d(z)\bigl(a(z)^N-b(z)^N\bigr)
       =d(z)\bigl(a(z)-b(z)\bigr)^N\ne0 .
\]
The tower extension contains the $\ell^m$th powers of the old Galois
group in the image of restriction for some $m$. This proves
$E_{-\alpha}(q_3)$ and $E_{-\sigma}(q_2)$ from the particular witnesses
constructed in RFCR1--2.
% END RFCR3

% BEGIN RFCR4
\paragraph{RFCR4. The seeded four-link chain on every later deep tower.}
Take any later deep $T$. Apply the seed promise of RFCR1 at this $T$,
with the original prime $w$: it lies in $W_2$ and hence dies on $T$.
This gives a finite coefficient extension and one integral seed pair
at $\iota$. Write $z$ for its reduction into the original $k$. Its
residual pair is obtained by reducing the integral pair and composing
the finite residue-field embedding into $k$. The coefficients agree
by compatibility of the original ambient maps; reduction mod the
maximal ideal gives local constancy. Thus $z$ satisfies the residual
conditions on the same empty frame. It lies above $s$ and satisfies
$\mathrm{InA}$, $E_{-\beta}$ and $E_{-\tau}$. Here
$\mathrm{InA}$ means representation support
$\{-\beta,-\sigma,-\tau\}$ and a nonzero $10$ coordinate of every
local flag. This same $z$ works for all the permitted residual $u$.

Take such a $u$. The initial global lattice shape and the prescribed
flag equations give either
\[
 (E_{-\sigma}(u),E_\beta(u))\quad\text{or}\quad
 (\mathrm{InV}(u), E_{-\alpha}(u)).
\]
In the second alternative $\mathrm{InV}$ means representation support
$\{\beta,-\alpha,-\tau\}$ with flags supported on $Q_0$ and all
$n_{20}\ne0$. To recall the restriction argument of UX2, a normalized
single-root cocycle that vanishes on the tower subgroup also vanishes
on conjugates of its separating element. The cocycle identity then
forces it to vanish globally, since the two diagonal values there
differ. In the second lattice alternative the flag equation in (UX.13)
propagates $n_{20}\ne0$. If the restricted $10$ entry vanished
identically, the global $10$ entry would vanish by the same argument;
the other first-column identity would make $n_{10}/n_{20}$ solve the
local equation excluded in the global lattice shape. Indeed on the initial
decomposition group at $\ell$, the $20$ flag equation reads
$n_{20}\bar\chi_0=r_{20}+\bar\chi_2 n_{20}$. Here $r_{20}=0$
and $n_{20}\ne0$, so $\bar\chi_0=\bar\chi_2$. This proves the
stated alternatives for the original $u$.

By RFCR3 this $u$ is $q_0$. The successive matching conditions are
\[
 u\ \xrightarrow{\mathrm{SamePs}}\ q_3\
   \xrightarrow{\mathrm{SameAl}}\ q_4\
   \xrightarrow{\mathrm{SamePs}}\ q_2\
   \xrightarrow{\mathrm{SameAl}}\ z .
\]
The last equality uses that $q_2$ restricts $p_2$, $z$ lies above $s$,
and RFCR1 preserves the $10$ data of $s$. We check goodness of every
endpoint. For $N_P$ support, $E_{-\sigma}$ is detected by the $20$
entry: the multiplier equation gives zero in $31$ whenever $20=0$.
For $N_M$ support the corresponding test for $E_{-\alpha}$ is the
$10$ entry, using $32$ as its reflected partner.
\begin{itemize}
\item At $u$ the alternatives just proved each imply ordinary goodness.
At $q_3$ the $-\alpha$ witness survives by RFCR3. If $u$ has
$E_{-\sigma}$, the first match transfers it to $q_3$. Otherwise
the vanishing of $20$ in the representation and nonvanishing of all
$20$ flags transfer $\mathrm{InV}$ to $q_3$, using the reflected
$31$ vanishing. Both alternatives give ordinary goodness.
\item At $q_4$ the middle $10$ match transfers $E_{-\alpha}$ from
$q_3$, and the subsequent $20$ match transfers $E_{-\sigma}$ from
$q_2$. Their conjunction gives ordinary goodness.
\item At $q_2$, $\mathrm{SameAl}$ with $z$ kills the $10$ and hence
the $32$ representation entries. Removing these from its $Q_0$
support leaves support in the required set for $\mathrm{InA}$;
its $10$ flags are nonzero by the same match. With its
$E_{-\sigma}$ this is the second $\mathrm{Good4}$ clause.
The last point $z$ has the third clause by its seed properties.
\end{itemize}
For each $\mathrm{SamePs}$ match the two pairs lie in $N_P$ with
flags in $Q_0$; for each $\mathrm{SameAl}$ match both representations
have $N_M$ support. Use the corresponding actual polynomial lines
of UX4. For clarity, interpolate the representation and flag matrices
entrywise with parameter $t$. The two endpoints on a $P$ line have
identical diagonal blocks in the ordering $(0,2\mid 1,3)$ (up to
the choice of which block comes first), and their differences occupy
one off-diagonal block. For an $M$ line use the ordering
$(0,1\mid2,3)$; matching $10$ determines the other diagonal block by
the multiplier equation. The differences of the lower flags also
lie in the relevant square-zero block, since they are lower
symplectic unipotents and have matching $20$ or $10$ coordinates.
Products and adjoints therefore give the interpolated endpoint
identities as in (UX.22). In particular conjugation by the interpolated
flag has the prescribed upper shape and diagonal characters. The other
frame conditions follow by evaluating the fixed entries and the
locally constant coefficient functions. This gives four links over
$k[t]$ in the order displayed, with every destination good and the
starting point good. Concatenation gives the literal $\mathrm{Chain4}$
from $u$ to the reduction of the one chosen seed. The construction
of $T_1$ preceded every later $T$, and the seed on $T$ preceded $u$.
This proves $\mathrm{RankFourChainReductionDense}$ under its source
cohomology, fitness and seed hypotheses.
% END RFCR4
